\section{Comparación con EP139: del marco de historia técnica a la acción del laboratorio}

La entrevista con Su Yu (苏煜) en EP139 ofrece una historia técnica de los Agent: cómo el Agent, entendido como un sistema con entidad, entorno y acción orientada a objetivos, pasó de logical agent y semantic parsing a language agent y universal digital agent. La entrevista con Luo Fuli (罗福莉) en EP138 responde otra pregunta: cuando este paradigma llega al interior de un laboratorio de modelos, cómo debe actuar el equipo.

\noindent\begin{tabularx}{\textwidth}{p{0.22\textwidth}p{0.34\textwidth}X}
\toprule
Dimensión & EP139 Su Yu & EP138 Luo Fuli \\
\midrule
Nivel de atención & Genealogía técnica y marco conceptual de los Agent & Recursos, organización y línea de modelos de un AI Lab. \\
Giro central & El OpenClaw Moment forma el consenso sobre los Agent & OpenClaw cambia el paradigma de investigación del equipo. \\
Mecanismos clave & Language Agent, disolución de fronteras, continual learning & Post-train, RL infra, orquestación de varios modelos, asignación de GPU. \\
Problemas de implementación & reliability, cost, deployment, world model & Tasa de finalización de extremo a extremo, costo y velocidad, igualdad organizativa, retroalimentación del entorno. \\
\bottomrule
\end{tabularx}

\begin{knowledgebox}{Cómo leerlas en conjunto}
Primero lee EP139 para entender por qué el Agent es un problema técnico de largo plazo; después lee EP138 para entender cómo un equipo de modelos ajusta sus recursos, su sistema de entrenamiento y su estructura organizativa una vez que ese problema técnico se convierte en consenso de la industria.
\end{knowledgebox}

\subsection{Resumen de la sección}

Los dos episodios forman un par: EP139 es el mapa conceptual de los Agent y EP138 es el mapa de acción de un laboratorio de modelos en la era de los Agent. La cola de Zhang Xiaojun AI puede seguir organizándose a partir de estos dos mapas.
